\documentclass[11pt]{article}
\usepackage{amsmath,amssymb}
\title{SF boundary insertions and original Wilson normalization}\author{}\date{}
\begin{document}\maketitle
\paragraph{Primary conventions checked.}
Sources: Luescher--Narayanan--Weisz--Wolff,
\texttt{https://arxiv.org/pdf/hep-lat/9207009}, sections3.4--3.5,4.5;
and Bode--Weisz--Wolff,
\texttt{https://arxiv.org/pdf/hep-lat/9809175}, sections2--3.
The former distinguishes ultraviolet subtraction from boundary
$O(a)$ improvement, with spatial boundary weight $c_s/2$ and attached
temporal weight $c_t$. The latter uses periodic space, time-Dirichlet
boundaries, and a special constant Abelian SU(3) family, computing
coupling coefficients using finite-lattice diagrams and extrapolation.
These statements do not supply an all-relative-cube norm or pressure.

\paragraph{Raw operators and ratio.}
Use the papers' oriented plaquette trace convention. Let $S_W/g_0^2$
be the base SF Wilson action, $O_t$ the literal sum over attached
temporal plaquettes, and $O_s$ HALF the sum over spatial boundary
plaquettes. Then
\[
 S_c=g_0^{-2}[S_W+(c_t-1)O_t+(c_s-1)O_s],\qquad
 \frac{Z_c}{Z_W}=E_W e^{-g_0^{-2}[(c_t-1)O_t+(c_s-1)O_s]}. \tag{1}
\]
No scalar or conditional denominator is removed. $O_s$ depends only
on prescribed boundary links. It vanishes identically on the special
commuting constant boundary family of9809175, but not for general
records. Changing the SF half allocation to another original Wilson
block allocation has an explicit boundary-dependent scalar too.

\paragraph{Two-loop insertion bookkeeping.}
Treat $c=c_t$ first as an independent parameter and write
\[
 \Gamma(g_0,c)=g_0^{-2}\Gamma_0(c)+\Gamma_1(c)
                               +g_0^2\Gamma_2(c)+O(g_0^4).
\]
Substitute $c=1+c_1g_0^2+c_2g_0^4+\cdots$. Taylor expansion yields
\[
 \Gamma_1^{\rm imp}=\Gamma_1^W+c_1\Gamma_{0,c},
\]
\[
 \boxed{\Gamma_2^{\rm imp}=\Gamma_2^W+c_1\Gamma_{1,c}
       +c_2\Gamma_{0,c}+\tfrac12c_1^2\Gamma_{0,cc}.} \tag{2}
\]
All derivatives are at $c=1$, with the actual minimum retained.
Let $A$ be its interior Hessian on the true quotient and
$j=\nabla_I O_t(U_*)$. Since the action is linear in $c$,
\[
 \Gamma_{0,c}=O_t(U_*),\qquad \Gamma_{0,cc}=-j^*A^{-1}j. \tag{3}
\]
The last term is the classical minimum displacement. It need not
vanish when the original $c=1$ plaquettes are constant.
$\Gamma_{1,c}$ includes moving-minimum, determinant, ghost and measure
response. Equivalently, in the ORIGINAL Wilson law,
$E_WO_t=O_t(U_*)+g_0^2\Gamma_{1,c}+\cdots$ and
$\operatorname{Var}_WO_t=g_0^2j^*A^{-1}j+\cdots$.
The connected expansion of (1) reproduces the negative quadratic
insertion in (2), independently of coupling inversion.

For several weights retain the full mixed Hessian
$\frac12\sum_{r,s}c_{r,1}c_{s,1}\Gamma_{0,rs}$.
The fixed spatial operator has zero interior gradient, hence zero
mixed variance/minimum response, but its scalar and source derivatives
remain. Returning to original weights requires ALL insertion terms
in (2), not merely removal of the $c_2$ term.

\paragraph{Free energy versus inverted coupling.}
At the specified parameter value let $k=\Gamma_0'$ and
$\bar g^2=k/\Gamma'$. If
$\bar g^2=g_0^2+m_1g_0^4+m_2g_0^6+\cdots$, then
\[
 m_1=-\Gamma_1'/k,\qquad
 \boxed{m_2-m_1^2=-\Gamma_2'/k.} \tag{4}
\]
Thus $m_2$ is not the free-energy two-loop coefficient. The primary
equation2.9 explicitly separates
$m_2-m_1^2=m_2^a+c_1m_2^b+c_1^2m_2^c+c_2m_2^d$.
These four terms are respectively
$-\Gamma_2^{W\prime}/k$, $-\Gamma_{1,c}'/k$,
$-\Gamma_{0,cc}'/(2k)$ and $-\Gamma_{0,c}'/k$.
The nonzero $c_1^2m_2^c$ is a genuine quadratic boundary insertion,
not the separate $m_1^2$ coupling inversion term.
The reported $m_1^b=m_2^d=-2a/L$ verifies the tree boundary count.
The table's $m_2^a$ isolates the unimproved contribution on that
particular family; an already improved $m_2$ does not.

This computes a specified first parameter derivative. It is neither
the arbitrary second boundary Hessian needed here nor the absolute
coefficient: the source-independent integration constant is lost.
Such a constant cannot be reconstructed from the running coupling,
and remains necessary in an absolute S94 comparison.

\paragraph{Transfer limit.}
The formulas above are valid insertion bookkeeping on the SAME
finite geometry/domain. They do not transfer periodic-spatial SF
asymptotics to a complete relative cube or prove uniformity over
boundary directions and histories. Boundary allocation, bulk coupling
conversion, original Wilson normalizers and every absolute scalar
remain to be compared. No growing-$R$ or nonlinear estimate is claimed.
\paragraph{A boundary-norm warning.}
On the complete
relative cube, put a unit Cartan tangent on the single boundary link
$e=((0,1,1,1),1)$ and zero on all other boundary links.
Four fixed $x_0=0$ face plaquettes (two in plane12 and two in plane13)
have unit squared flux. The zero-interior extension adds just one
inward $01$ plaquette, for total energy $5$. Thus the true harmonic energy
lies in $[4,5]$, and the fixed-face energy divided by the total is
at least $4/5$, uniformly in $R\ge4$.
Therefore the smooth-source $O(a/L)$ boundary improvement scale is
not uniform smallness in the complete energy tangent norm. This does
not contradict SF asymptotics or a uniform bounded boundary form;
it forbids silently deleting $c_s$ in a general-domain comparison.
\end{document}
